\begin{pastexam}[label={px:df3-096-ds-05}]{地方三等資訊處理 96 年 第五題}
（申論；20 分）

有一遞迴函數如下：
\begin{minted}[fontsize=\small, style=perldoc]{c}
int g(int n)
{
  if (n < 2)
     return n;
  else
     return g(n-1) + g(n-2) + n;
}
\end{minted}
\begin{itemize}[leftmargin=2em, labelsep=0.3em, itemsep=2pt]
\item[㈠] 請寫出 \texttt{g(6)} 的傳回值。（3 分）
\item[㈡] 請算出此遞迴函數的時間複雜度（time complexity）。（7 分）
\item[㈢] 請寫一全等於上述遞迴函數之線性時間（linear time）函數。（10 分）
\end{itemize}
\end{pastexam}

\begin{pastsoln}
\textbf{㈠} 由小往上算，$g(0) = 0$、$g(1) = 1$，之後 $g(n) = g(n-1) + g(n-2) + n$：
\[
\begin{aligned}
g(2) &= g(1) + g(0) + 2 = 1 + 0 + 2 = 3, & g(3) &= 3 + 1 + 3 = 7, & g(4) &= 7 + 3 + 4 = 14,\\
g(5) &= 14 + 7 + 5 = 26, & g(6) &= 26 + 14 + 6 = \mathbf{46}.
\end{aligned}
\]
\texttt{g(6)} 傳回 $46$。

\textbf{㈡} 每次呼叫除了遞迴之外只做常數時間的工作，所以時間和「呼叫次數」成正比。 設
$C(n)$ 是算 \texttt{g(n)} 總共的呼叫次數（含自己）：
\[
C(0) = C(1) = 1, \qquad C(n) = C(n-1) + C(n-2) + 1 \quad (n \ge 2).
\]
兩邊加 $1$：令 $D(n) = C(n) + 1$，得 $D(n) = D(n-1) + D(n-2)$，$D(0) = D(1) = 2$。 這是費氏
數列的兩倍：$D(n) = 2F_{n+1}$（$F_1 = F_2 = 1$），所以
\[
C(n) = 2F_{n+1} - 1, \qquad C(0..6) = 1, 1, 3, 5, 9, 15, 25 .
\]
費氏數的特徵方程式 $x^2 = x + 1$ 的根是 $\varphi = (1 + \sqrt5)/2 \approx 1.618$ 和
$\hat\varphi = (1 - \sqrt5)/2$，$|\hat\varphi| < 1$，所以 $F_n = (\varphi^n - \hat\varphi^n)/\sqrt5
= \Theta(\varphi^n)$（資結書第~2.5~節的線性遞迴）。 因此時間
\[
T(n) = \Theta(C(n)) = \Theta(\varphi^n) \approx \Theta(1.618^n),
\]
是指數時間。 不想解特徵方程式也可以夾出範圍：$C(n-1) \ge C(n-2)$，所以
$2C(n-2) \le C(n) - 1 \le 2C(n-1)$，得 $C(n) = \Omega(2^{n/2})$ 且 $C(n) = O(2^n)$，一樣是指數。
慢的原因是同一個 \texttt{g(k)} 被重算很多次：算 \texttt{g(6)} 時 \texttt{g(2)} 被呼叫了
$F_5 = 5$ 次。

\textbf{㈢} 由小往上算，只記住最近的兩個值（和資結書第~18.1~節「由下而上」的動態規劃
同一個想法）：
\begin{minted}[fontsize=\small, style=perldoc]{c}
int g_linear(int n)
{
    if (n < 2)              /* 和原函數相同：n < 2 時傳回 n（含負數） */
        return n;
    int prev = 0;           /* g(i-2)，一開始是 g(0) */
    int cur  = 1;           /* g(i-1)，一開始是 g(1) */
    for (int i = 2; i <= n; i++) {
        int next = cur + prev + i;   /* g(i) = g(i-1) + g(i-2) + i */
        prev = cur;
        cur  = next;
    }
    return cur;             /* 迴圈結束時 cur = g(n) */
}
\end{minted}
\textbf{為什麼全等}：迴圈不變量是「第 $i$ 次迴圈開始時 \texttt{prev} $= g(i-2)$、\texttt{cur}
$= g(i-1)$」。 一開始 $i = 2$，\texttt{prev} $= g(0) = 0$、\texttt{cur} $= g(1) = 1$ 成立；
迴圈本體算出 $g(i)$ 後左移一格，$i$ 加一，不變量仍然成立。 結束時 $i = n + 1$，\texttt{cur}
$= g(n)$。 $n < 2$ 的情況直接照原函數傳回 $n$。 所以對每個 $n$ 傳回值都和原函數相同。
以 $n = 6$ 追蹤：$(\texttt{prev}, \texttt{cur})$ 依序是 $(0,1) \to (1,3) \to (3,7) \to (7,14)
\to (14,26) \to (26,46)$，傳回 $46$，和㈠一致。

\textbf{時間}：迴圈 $n - 1$ 次，每次常數時間，$\Theta(n)$；空間只用三個變數，$O(1)$。

補充：令 $h(n) = g(n) + n + 3$，代入可得 $h(n) = h(n-1) + h(n-2)$，$h(0) = 3 = F_4$、
$h(1) = 5 = F_5$，所以 $g(n) = F_{n+4} - n - 3$（$n \ge 0$）。 例如 $g(6) = F_{10} - 9 = 55 - 9
= 46$。 這也說明 $g(n)$ 本身是指數成長，$n \ge 43$ 時就超出 32 位元 \texttt{int} 的範圍；
在那之前兩個函數的傳回值完全相同。
\end{pastsoln}
